\documentclass[a4paper,12pt]{article}

\usepackage[left=1.4cm,top=1.2cm,bottom=1.5cm,right=1.2cm]{geometry}%


\usepackage[latin1]{inputenc}
\usepackage{amsmath}
%\usepackage{babel}
\topmargin 0truecm
\textheight 23truecm
\textwidth 16.5truecm
\oddsidemargin 0truecm
\parskip 1truemm
%\parsep 1truemm
\parindent 0em
\usepackage{amssymb}

\def\d{\,{\rm d}}

\def\RR{{\mathbb R}}
\def\CC{{\mathbb C}}
\def\NN{{\mathbb N}}
\def\ZZ{{\mathbb Z}}
\def\EE{{\mathbb E}}
\newcounter{problema}
\newcommand{\prob}[1]{\vspace{0.33cm}\stepcounter{problema}
                 \noindent{\bf Problema \arabic{problema}:}{\it #1}}

\begin{document}
\renewcommand{\labelenumi}{\bf \alph{enumi})}
\renewcommand{\labelenumii}{\alph{enumi}$_\arabic{enumii}$)}
\thispagestyle{empty}
\begin{center}
{\large {\bf M\'etodos Num\'ericos}}\\



\vspace*{1.5cm} 
\rule{0em}{2.em}{\bf Gu\'{i}a 0} \\ Marzo de 2026
\end{center}
\vspace*{.5cm} 

\vspace*{.5cm} 
\textbf{NOTA:} Los ejercicios de esta gu\'ia no se resuelven programando, sino con l\'apiz y papel.
\vspace*{.5cm} 



\prob{} Elabore un algoritmo que permita ingresar el n\'umero de partidos
ganados, perdidos y empatados de un equipo arbitrario de f\'utbol, y que
devuelva los puntos obtenidos (en cada partido un equipo obtiene $3$ puntos
si gana, $1$ punto si empata y ning\'un punto si pierde).


\prob{} Elabore un algoritmo que lea $3$ n\'umeros enteros diferentes entre
s\'{\i} y determine cu\'al de ellos es el mayor.


\prob{} Elabore un algoritmo que lea un n\'umero de dos cifras en su 
representaci\'on decimal y que devuelva el n\'umero de unidades y de decenas
que lo componen.  



\prob{} Escriba, con papel y l\'apiz, un algoritmo que le permita 
transformar un n\'umero entero de la base decimal a la binaria. 
Incluya la entrada, la salida y el criterio de
parada.

\prob{} Escriba, con papel y l\'apiz, un algoritmo que le permita transformar
la base binaria a la decimal.  Incluya la entrada, la salida y el criterio de
parada.


\prob{} Escriba los siguientes n\'umeros en base binaria:


\begin{itemize}

\item[a)] $(29)_{10}$
 
\item[b)] $(276)_{10}$
 
\item[c)] $(0.465)_{10}$

\item[d)] $(26.75)_{10}$

\end{itemize}



\prob{} Escriba los siguientes n\'umeros en base decimal:


\begin{itemize}

\item[a)] $(1001101)_{2}$
 
\item[b)] $(11)_{2}$
 
\item[c)] $(0.1011)_{2}$

\item[d)] $(1010011.10110101)_{2}$

\end{itemize}




\prob{} Escriba, en papel y con l\'apiz, los siguentes n\'umeros enteros en base binaria:
\begin{enumerate}
\item $-16$ en complemento de $2$ con $6$ bits
\item $+13$ en signo y magnitud con $5$ bits,
\item $-64$ en complemento de dos con $7$ bits,
\item $+12$ en signo y magnitud con $6$ bits.
\item $-29$ en signo y magnitud con $6$ bits.
\item $-29$ en complemento de dos con $6$ bits.
\end{enumerate}



\prob{}   Evaluar las siguientes operaciones
usando matem\'atica exacta y {\em matem\'atica  de enteros de 1 byte} (complemento de dos). 

\begin{enumerate}
\item {\tt 127 + 1}
\item {\tt 100*2}  
\item {\tt -100 - 40}  
\end{enumerate}% 2



\prob{} Pasar los siguientes n\'umeros a punto flotante en base 2, simple precisi\'on, 
seg\'un la norma $IEEE754$:

\begin{itemize}

\item[a)] $(0.125)_{10}$
 
\item[b)] $(-12.75)_{10}$
 
\item[c)] $(2890.2)_{10}$
 
\end{itemize}


\prob{} Calcular los errores absolutos y relativos de aproximar $p$ por $p^*$.

\begin{itemize}

\item[a)] $p =0.3000 \times 10^1$ ; $p^* =0.3100 \times 10^1$
 
\item[b)] $p =0.3000 \times 10^{-3}$ ; $p^* =0.3100 \times 10^{-3}$
 
\item[c)] $p =0.3000 \times 10^4$ ; $p^* =0.3100 \times 10^4$

\item[d)] $p = \pi$ ; $p^* = 22/7$
 
\item[e)] $p = \pi$ ; $p^* = 3.1416$
 
\item[f)] $p = e$ ; $p^* = 2.718$
 
\item[g)] $p = \sqrt{2}$ ; $p^* = 1.414$
 
\end{itemize}




\prob{} Realice los siguientes c\'alculos:

\begin{itemize}

\item[i)] exactamente
 
\item[ii)] usando aritm\'etica de corte de 3 d\'igitos
 
\item[iii)] usando aritm\'etica de redondeo de 3 d\'igitos
 
\item[iv)] Calcule los errores relativos en las partes $(ii)$ y $(iii)$
 
\end{itemize}

\begin{itemize}

\item[a)] $\frac{4}{5} + \frac{1}{3}$
 
\item[b)] $\frac{4}{5} \times \frac{1}{3}$
 
\item[c)] $(\frac{1}{3} - \frac{3}{11}) + \frac{3}{20}$
  
\end{itemize}




% p
\prob{} Use aritm\'etica de 4 d\'igitos (redondeando) para simular el problema del c\'alculo computacional
de $\pi - \frac{22}{7}$. 
Luego calcule el error absoluto y el error relativo de la representaci\'on de $\pi$ y de $\frac{22}{7}$,
y el error relativo de la diferencia.



\prob{}  Se dispone de dos f\'ormulas para obtener el valor $x^*$ de la 
intersecci\'on de la recta con el eje {\em x}:
\[
x^* = \frac{x_0y_1 - x_1 y_0}{y_1-y_0} \qquad \qquad y  \qquad \qquad
x^* = x_0 - \frac{(x_1-x_0) y_0}{y_1-y_0}
\]
\begin{enumerate}
\item Muestre que ambas f\'ormulas son algebraicamente correctas
\item Use los datos $(x_0,y_0) = (1.31,3.24)$ y $(x_1,y_1) = (1.93,4.76)$ y aritm\'etica de tres d\'igitos
para calcular $x^*$ de ambas formas.  Cu\'al m\'etodo es mejor? Por qu\'e'? 
\end{enumerate}








\end{document}


